\documentclass[10pt,a4paper]{amsart}
\usepackage{amsmath,amssymb,amsthm,graphicx}
% Retained preamble for the Stacks Project chapter {CH} (shared project preamble).
% The chapter file pinned-original-source.tex only carries \input{{preamble}} on its first line,
% so no preamble text lives inside the chapter itself.  The Stacks Project uses one single shared
% preamble for every chapter: the snapshot stages it three times (sources/stacks-chow/preamble.tex,
% sources/stacks-etale-cohomology/preamble.tex and focused-seeds/AGP/sources/stacks-snapshot-20261011/
% preamble.tex) and those copies are identical in size.  The block below is copied mechanically,
% line for line, from sources/stacks-chow/preamble.tex lines 14-18 and 154-217: the Xy-pic setup and
% the native theorem declarations together with the mathematical macros.  Nothing was added, removed
% or redefined; the wrapper class and the amsmath/amssymb/amsthm/graphicx base packages are supplied
% by the assembler.
% For commutative diagrams we use Xy-pic
% Chapter: stacks-cohomology.tex -- Cohomology of Sheaves.
\usepackage[all]{xy}

% We use 2cell for 2-commutative diagrams.
\xyoption{2cell}

% Theorem environments.
%
\theoremstyle{plain}
\newtheorem{theorem}[subsection]{Theorem}
\newtheorem{proposition}[subsection]{Proposition}
\newtheorem{lemma}[subsection]{Lemma}

\theoremstyle{definition}
\newtheorem{definition}[subsection]{Definition}
\newtheorem{example}[subsection]{Example}
\newtheorem{exercise}[subsection]{Exercise}
\newtheorem{situation}[subsection]{Situation}

\theoremstyle{remark}
\newtheorem{remark}[subsection]{Remark}
\newtheorem{remarks}[subsection]{Remarks}

\numberwithin{equation}{subsection}

% Macros
%
\def\lim{\mathop{\mathrm{lim}}\nolimits}
\def\colim{\mathop{\mathrm{colim}}\nolimits}
\def\Spec{\mathop{\mathrm{Spec}}}
\def\Hom{\mathop{\mathrm{Hom}}\nolimits}
\def\Ext{\mathop{\mathrm{Ext}}\nolimits}
\def\SheafHom{\mathop{\mathcal{H}\!\mathit{om}}\nolimits}
\def\SheafExt{\mathop{\mathcal{E}\!\mathit{xt}}\nolimits}
\def\Sch{\mathit{Sch}}
\def\Mor{\mathop{\mathrm{Mor}}\nolimits}
\def\Ob{\mathop{\mathrm{Ob}}\nolimits}
\def\Sh{\mathop{\mathit{Sh}}\nolimits}
\def\NL{\mathop{N\!L}\nolimits}
\def\CH{\mathop{\mathrm{CH}}\nolimits}
\def\proetale{{pro\text{-}\acute{e}tale}}
\def\etale{{\acute{e}tale}}
\def\QCoh{\mathit{QCoh}}
\def\Ker{\mathop{\mathrm{Ker}}}
\def\Im{\mathop{\mathrm{Im}}}
\def\Coker{\mathop{\mathrm{Coker}}}
\def\Coim{\mathop{\mathrm{Coim}}}

% Boxtimes
%
\DeclareMathSymbol{\boxtimes}{\mathbin}{AMSa}{"02}

%
% Macros for moduli stacks/spaces
%
\def\QCohstack{\mathcal{QC}\!\mathit{oh}}
\def\Cohstack{\mathcal{C}\!\mathit{oh}}
\def\Spacesstack{\mathcal{S}\!\mathit{paces}}
\def\Quotfunctor{\mathrm{Quot}}
\def\Hilbfunctor{\mathrm{Hilb}}
\def\Curvesstack{\mathcal{C}\!\mathit{urves}}
\def\Polarizedstack{\mathcal{P}\!\mathit{olarized}}
\def\Complexesstack{\mathcal{C}\!\mathit{omplexes}}
% \Pic is the operator that assigns to X its picard group, usage \Pic(X)
% \Picardstack_{X/B} denotes the Picard stack of X over B
% \Picardfunctor_{X/B} denotes the Picard functor of X over B
\def\Pic{\mathop{\mathrm{Pic}}\nolimits}
\def\Picardstack{\mathcal{P}\!\mathit{ic}}
\def\Picardfunctor{\mathrm{Pic}}
\def\Deformationcategory{\mathcal{D}\!\mathit{ef}}

% --- appended 2026-10-12 (Chief): the project own definitions of the reference/slogan/history
% display environments, copied mechanically from the shared Stacks preamble lines 8-12
% (\usepackage{verbatim} plus the three \newenvironment lines, which the chapter preamble keeps
% verbatim from the project preamble). Nothing added, removed or redefined: these environments
% carry no mathematical content and the project itself typesets them as comments. Without these
% three lines a counted statement whose own text uses them fails with Environment slogan undefined.
\usepackage{verbatim}
\newenvironment{reference}{\comment}{\endcomment}
\newenvironment{slogan}{\comment}{\endcomment}
\newenvironment{history}{\comment}{\endcomment}
\usepackage{cleveref}
\crefname{originallemma}{original lemma}{original lemmas}
\Crefname{originallemma}{Original lemma}{Original lemmas}
\makeatletter
\newlabel{lemma-solution-in-finite-case}{{0D6A}{source}{}{sourcepointer-1}{}}
\newlabel{lemma-solution-in-finite-case@cref}{{[originallemma][0][]0D6A}{[1][0][]source}{}{}{}}
\makeatother
\title{Glueing along a well-ordered increasing union of opens}
\author{The Stacks Project Authors}
\date{stacks-project:cohomology.tex:local-snapshot-20261011; GNU FDL 1.2 or later}
\begin{document}
\maketitle
\noindent\textit{Editorial note.} All delivered mathematical text is the author's own, copied verbatim from the pinned Stacks Project chapter file cohomology.tex; no statement or proof step was rewritten, reordered or removed.\par\medskip
\noindent\textit{Editorial note.} The counted span cites the lemmas recorded in assembly.external\_references with canonical tags as navigation pointers.\par\medskip
\noindent\textit{Allowed standard prerequisites (editorial).}\begin{itemize}
\item Uniqueness of isomorphisms of solutions. Cohomology of Sheaves, Lemma cohomology-lemma-solution-in-finite-case. Two solutions of a system with vanishing negative exts are uniquely isomorphic.
\item Filtered colimits are exact. Modules, Section modules-section-colimits. Stalk computation of the colimit complex.
\end{itemize}\medskip
\ifdefined\hypertarget\hypertarget{sourcepointer-1}{}\fi
\noindent\textit{Source reference.} Original lemma 0D6A: sources/stacks-chow/tags.txt, source lines 16515--16515. Cross-chapter standard prerequisite cited by the author inside the counted span; it is used as a pointer and its own proof is not part of this item.\par\medskip
\begin{lemma}
\label{lemma-glueing-increasing-union}
Let $X$ be a ringed space. Let $E$ be a well ordered set and let
$$
X = \bigcup\nolimits_{\alpha \in E} W_\alpha
$$
be an open covering with $W_\alpha \subset W_{\alpha + 1}$
and $W_\alpha = \bigcup_{\beta < \alpha} W_\beta$ if $\alpha$ is not
a successor. Let $K_\alpha$ be an object of $D(\mathcal{O}_{W_\alpha})$
with $\Ext^i(K_\alpha, K_\alpha) = 0$ for $i < 0$.
Assume given isomorphisms
$\rho_\beta^\alpha :  K_\alpha|_{W_\beta} \to K_\beta$ in
$D(\mathcal{O}_{W_\beta})$ for all $\beta < \alpha$ with
$\rho_\gamma^\alpha = \rho_\gamma^\beta \circ \rho^\alpha_\beta|_{W_\gamma}$
for $\gamma < \beta < \alpha$.
Then there exists an object
$K$ in $D(\mathcal{O}_X)$ and isomorphisms
$K|_{W_\alpha} \to K_\alpha$ for $\alpha \in E$
compatible with the isomorphisms $\rho_\beta^\alpha$.
\end{lemma}

\begin{proof}
In this proof $\alpha, \beta, \gamma, \ldots$ represent elements of $E$.
Choose a K-injective complex
$I_\alpha^\bullet$ on $W_\alpha$ representing $K_\alpha$.
For $\beta < \alpha$ denote $j_{\beta, \alpha} : W_\beta \to W_\alpha$
the inclusion morphism. Using transfinite recursion we will construct for all
$\beta < \alpha$ a map of complexes
$$
\tau_{\beta, \alpha} :
(j_{\beta, \alpha})_!I_\beta^\bullet
\longrightarrow
I_\alpha^\bullet
$$
representing the adjoint to the inverse of the isomorphism
$\rho^\alpha_\beta : K_\alpha|_{W_\beta} \to K_\beta$.
Moreover, we will do this in such that for
$\gamma < \beta < \alpha$ we have
$$
\tau_{\gamma, \alpha} = \tau_{\beta, \alpha} \circ
(j_{\beta, \alpha})_!\tau_{\gamma, \beta}
$$
as maps of complexes. Namely, suppose already given $\tau_{\gamma, \beta}$
composing correctly for all $\gamma < \beta < \alpha$.
If $\alpha = \alpha' + 1$ is a successor, then we choose any map of complexes
$$
(j_{\alpha', \alpha})_!I_{\alpha'}^\bullet \to I_\alpha^\bullet
$$
which is adjoint to the inverse of the isomorphism
$\rho^\alpha_{\alpha'} : K_\alpha|_{W_{\alpha'}} \to K_{\alpha'}$
(possible because $I_\alpha^\bullet$ is K-injective)
and for any $\beta < \alpha'$ we set
$$
\tau_{\beta, \alpha} = \tau_{\alpha', \alpha} \circ
(j_{\alpha', \alpha})_!\tau_{\beta, \alpha'}
$$
If $\alpha$ is not a successor, then
we can consider the complex on $W_\alpha$ given by
$$
C^\bullet = \colim_{\beta < \alpha} (j_{\beta, \alpha})_!I_\beta^\bullet
$$
(termwise colimit) where the transition maps of the sequence
are given by the maps $\tau_{\beta', \beta}$ for
$\beta' < \beta < \alpha$. We claim that $C^\bullet$
represents $K_\alpha$. Namely, for $\beta < \alpha$ the restriction
of the coprojection $(j_{\beta, \alpha})_!I_\beta^\bullet \to C^\bullet$
gives a map
$$
\sigma_\beta : I_\beta^\bullet \longrightarrow C^\bullet|_{W_\beta}
$$
which is a quasi-isomorphism: if $x \in W_\beta$ then looking
at stalks we get
$$
(C^\bullet)_x =
\colim_{\beta' < \alpha}
\left((j_{\beta', \alpha})_!I_{\beta'}^\bullet\right)_x =
\colim_{\beta \leq \beta' < \alpha} (I_{\beta'}^\bullet)_x
\longleftarrow
(I_\beta^\bullet)_x
$$
which is a quasi-isomorphism. Here we used that taking stalks
commutes with colimits, that filtered colimits are exact, and
that the maps $(I_\beta^\bullet)_x \to (I_{\beta'}^\bullet)_x$
are quasi-isomorphisms for $\beta \leq \beta' < \alpha$.
Hence $(C^\bullet, \sigma_\beta^{-1})$ is a solution to the
system $(\{K_\beta\}_{\beta < \alpha},
\{\rho^\beta_{\beta'}\}_{\beta' < \beta < \alpha})$.
Since $(K_\alpha, \rho^\alpha_\beta)$ is another solution
we obtain a unique isomorphism $\sigma : K_\alpha \to C^\bullet$
in $D(\mathcal{O}_{W_\alpha})$ compatible with all our maps, see
Lemma \ref{lemma-solution-in-finite-case}
(this is where we use the vanishing of negative ext groups).
Choose a morphism $\tau : C^\bullet \to I_\alpha^\bullet$
of complexes representing $\sigma$. Then we set
$$
\tau_{\beta, \alpha} = \tau|_{W_\beta} \circ \sigma_\beta
$$
to get the desired maps. Finally, we take $K$ to be the object of the derived
category represented by the complex
$$
K^\bullet = \colim_{\alpha \in E} (W_\alpha \to X)_!I_\alpha^\bullet
$$
where the transition maps are given by our carefully constructed
maps $\tau_{\beta, \alpha}$ for $\beta < \alpha$.
Arguing exactly as above we see that for all $\alpha$
the restriction of the coprojection determines an isomorphism
$$
K|_{W_\alpha} \longrightarrow K_\alpha
$$
compatible with the given maps $\rho^\alpha_\beta$.
\end{proof}

\end{document}
